\chapter{Methodology}

\section{Development Approach}

This project is primarily centred on the design and development of a full stack software system. The main objective was to build and deliver a functional application, and the process was therefore guided by established software engineering practices. However, the development was supported by a focused research component, particularly in the areas of machine learning model selection and training, where decisions around algorithms, datasets, and evaluation strategies were informed by existing literature and exploratory experimentation before implementation. An empirical approach was applied throughout: rather than fixing a design upfront, each subsystem was developed, tested, measured, and refined based on observed performance.

The overarching development approach was agile, iterative, and incremental,
drawing on the principles of the Agile Manifesto \cite{agilemanifesto2001}
and lightweight practice suggested by Beck and Andres \cite{beck2004xp}.
Rather than attempting to design the entire system upfront and implement
it in sequence from start to finish, the project was broken into a series
of development cycles, each delivering a working vertical slice of
functionality that could be tested, reviewed, and built upon. This
approach was well suited to a project of this complexity, where the
requirements for the machine learning components could not be fully
specified in advance. The decision to use a multi label classifier for
topic tagging, for example, only emerged once the nature of the training
data was properly understood: the LeetCode corpus showed that problems
routinely carry several tags at once, which ruled out a single label
classifier from the outset. Similarly, the approach to difficulty labelling
only became clear after exploring the structure of the IBM CodeNet dataset
and noticing that while no explicit difficulty labels were available,
acceptance rate provided a usable proxy signal.

\section{Planning and Requirements}

Before any implementation began, a significant amount of upfront planning
was undertaken to establish the scope of the project and the technologies
that would be used. This planning work was produced entirely within Figma
and is documented in the \texttt{/docs/project\_planning} folder of the
project repository. The artefacts produced include a concept summary
outlining the high level vision for the platform, an original system
architecture diagram showing how the major components would relate to each
other, a planned tech stack document justifying the selection of each
major technology, prototype requirements describing the expected behaviour
of the system at a minimum viable product level, and a Docker container
plan specifying how the three services (database, backend, and frontend)
would be orchestrated. A Gantt chart was also produced, created in Canva,
mapping out the intended timeline for each phase of development from
initial research through to final testing and writeup.

This planning phase served as a storyboarding exercise for the entire
project. By working through the architecture visually before writing any
code, it was possible to identify dependencies between components early.
For instance, it became clear that the backend API contracts needed to be
defined before meaningful frontend development could begin, and that the
machine learning models needed to be trained and serialised before they
could be integrated into the inference pipeline. Identifying these ordering
constraints in advance meant that the sequence of sprints could be
arranged to minimise rework.

Requirements were determined through a combination of domain analysis and
personal specification. Since this was a solo project without an external
client, the requirements were derived by examining similar platforms,
identifying the gaps described in the introduction, and translating them
into concrete functional and non functional specifications. Functional
requirements included question generation via prompt, multiple choice and
coding challenge formats, topic tag prediction, difficulty prediction,
code execution with test case evaluation, user authentication, and a
social friends system. Non functional requirements included containerised
deployment for reproducibility, sub second machine learning inference
latency so that the enrichment step would not degrade the perceived speed
of question generation, and support for at least five programming languages
in the code runner so that the platform would serve a realistic spread of
coursework scenarios.

\section{Iterative Development Cycles}

Development proceeded through a series of loosely defined sprints, each
targeting a coherent piece of functionality. The early sprints focused on
establishing the foundational infrastructure: configuring the PostgreSQL
database, handling API calls to large language models, and building the
basic FastAPI skeleton with a working database connection. Once the
infrastructure was stable, subsequent cycles added the AI generation
pipeline, the machine learning training and integration work, the code
execution engine, the authentication system, and finally the social
features and frontend polish.

This incremental approach meant that the system was in a runnable and
testable state throughout most of the development period. At no point was
the entire codebase non functional for more than a short period while a
breaking change was integrated. This mattered because it allowed problems
to be caught early, when the context of what had just changed was still
fresh, rather than at the end of a long development phase when the source
of a bug could be difficult to trace. The principle at work here is
essentially the same one that underlies continuous integration: small,
frequent changes are cheaper to verify and cheaper to revert than large
ones.

Where technical problems arose, an empirical approach was taken. When the
topic tag classifier initially produced poor results, for example, the
response was not to switch to a different algorithm immediately but rather
to investigate the training data more carefully, examine the class
distribution, experiment with threshold tuning, and try different feature
engineering strategies before drawing conclusions. This reflects a
scientific problem solving habit even within a software development
context: observe, hypothesise, test, and iterate. In several cases the
root cause turned out to be in the data rather than the model. Rare topic
tags with only a handful of examples were producing unreliable predictions
not because the classifier was wrong but because there simply was not
enough signal to learn from. The eventual fix was a minimum sample
threshold filter on the tag set rather than a change in algorithm.

\section{GitHub Workflow}

GitHub was used throughout the project as both the version control system
and the primary project management tool. The repository followed a flexible
branching strategy centred around a stable \texttt{main} branch, with
additional branches created as needed to support ongoing development.
Rather than adhering to a strict two branch model, new branches were often
created to address specific features or areas of the system, such as social
functionality, or to work on tasks associated with particular GitHub
issues. This lightweight approach suited a solo project, where the
overhead of maintaining long lived parallel branches would have outweighed
the benefits.

Work was tracked using GitHub Issues, where each issue represented a
discrete unit of work, including new features, bug fixes, refactoring
tasks, or research activities. Issues were organised using a set of labels
tailored to the project, including categories such as classifier model,
bugs, coding sandbox, and MCQ quiz functionality, alongside default labels.
This labelling approach allowed for efficient filtering and provided a
clear overview of the types of work remaining at any stage of development.

Milestones were used to group related issues into time bounded delivery
targets aligned with the overall project plan. Each milestone corresponded
to a major phase of development, such as the MVP launch, MCQ system
completion, classifier integration, and final dissertation submission.
This structure provided a high level view of progress and helped ensure
that development remained aligned with key deadlines, including the
Christmas seminar presentation and the April submission.

When beginning work on a task, a new branch was typically created,
sometimes linked to a specific issue, and named in a way that reflected
the functionality being developed, for example \texttt{social\_features}
or \texttt{code\_sandbox\_v2}. Development was carried out through
incremental commits with descriptive messages. Upon completion, changes
were merged back into the main codebase. This workflow maintained a clear
and navigable project history while supporting parallel development across
different features, and it left a commit trail that serves as evidence of
sustained effort across the full development period rather than a last
minute sprint.

\section{Continuous Integration Pipeline}

A continuous integration pipeline was configured using GitHub Actions
\cite{githubactions2025} and is defined in the \texttt{.github/workflows}
directory of the repository. The pipeline triggers on every push to the \texttt{main} branch and on every pull request targeting it. This ensured that no code could be merged into the \texttt{main} branch without first passing automated checks.

The pipeline is divided into three parallel jobs. The frontend job
installs Node.js 22 and the project dependencies, runs ESLint to check
for code quality issues, performs a TypeScript type check using
\texttt{tsc --noEmit} to catch type errors without producing output files,
and then runs the production build using Vite \cite{vite2025} to confirm
that the application compiles successfully. Passing build artifacts are
uploaded for a retention period of seven days. The backend job sets up
Python 3.11, runs the Black formatter in check mode to verify that the
code conforms to the project's formatting standard, runs Flake8 to catch
syntax errors and undefined name references, and compiles every Python
file in the backend directory to validate syntax. The ML models job runs
a similar set of checks across the \texttt{ml\_models} directory. A final
summary job runs only if all three parallel jobs succeed, providing a
clear single status indicator for the overall health of the pipeline.

This setup served two important purposes. First, it enforced a baseline
level of code quality automatically, removing the need for manual checks
on every change and catching regressions immediately. Second, it provided
confidence when merging feature branches that the integration would not
break the build. The use of \texttt{continue-on-error: true} on the
formatting and complexity checks meant that these were informational
rather than blocking, allowing the pipeline to surface stylistic feedback
without preventing progress when minor formatting divergences were present.
This calibration matters: a pipeline that fails on every whitespace
difference ends up being ignored by its author, whereas a pipeline that
fails only on genuinely broken builds retains its value as a signal.

\section{Validation and Testing}

Testing was carried out primarily through manual integration testing and
through the empirical evaluation of the machine learning models on held
out test sets. Given that this was a solo project and the focus of the
evaluation chapter is on measurable outcomes rather than on test coverage
metrics, a full unit test suite was not implemented. A representative
sample of pytest tests does exist under \texttt{backend/tests}, covering
core logic such as JWT token round tripping and password hashing, but the
bulk of validation was done at the integration level, where the behaviour
that matters to the end user is actually observable. The continuous
integration pipeline's syntax validation and build checks served as a
lightweight automated safety net across all three components of the
codebase.

The machine learning models were evaluated rigorously using train,
validation, and test splits. For the topic tag classifier a sixty, twenty,
twenty split was used, with the validation set serving as the basis for
per label decision threshold optimisation before final evaluation on the
held out test set. The rationale for this approach is set out in detail
in the System Design chapter, but the key point is that tuning thresholds
on the test set would have leaked information and produced optimistic
estimates of generalisation. Holding out a separate validation set for
threshold selection, then reporting test set numbers only once at the
end, is the standard cross validation protocol described by Kohavi
\cite{kohavi1995crossval}. For the difficulty classifier a five fold cross
validation strategy was used during hyperparameter search via
\texttt{RandomizedSearchCV} \cite{bergstra2012random}, which is a more
computationally efficient alternative to exhaustive grid search when the
hyperparameter space is large.

The code execution engine was tested manually by submitting known correct
and intentionally incorrect solutions across each supported language and
verifying that the test case evaluation logic returned the expected
outcomes. The end to end question generation pipeline was tested by
issuing prompts and inspecting the enriched responses, checking that the
tags and difficulty labels returned were plausible for the generated
content. Where prompts produced questions that the model then tagged
implausibly, those cases were logged and fed back into a later round of
prompt engineering or threshold tuning, which over time reduced the rate
of surprising outputs.

\section{Technology Selection Criteria}

The choice of specific technologies was guided by a small number of
consistent criteria. First, each major technology had to be open source
or freely available at the project's scale, because a solo student project
cannot justify paid tooling beyond what is strictly necessary. Second,
each technology had to have active community support and documentation
sufficient to resolve problems without a commercial support contract.
Third, where a choice was made between two roughly equivalent options,
the option with stronger type safety, stronger validation, or stronger
reproducibility was preferred, on the grounds that these properties
compound favourably as a codebase grows. FastAPI over Flask, TypeScript
over plain JavaScript, and Pydantic validation on every request boundary
are all choices that reflect this preference.